\documentclass{article}
\usepackage[utf8]{inputenc}
\usepackage{amsmath}
\usepackage{geometry}
\geometry{margin=2cm}
\begin{document}

\section*{Análise da Questão 156 — ENEM 2024 — Matemática}

\section{Enunciado}
O uso de aplicativos de transporte tem sido uma alternativa à população que busca preços mais competitivos para se locomover, principalmente nas grandes cidades. As formas usadas para determinar o valor cobrado por cada viagem variam de um aplicativo para outro, mas, em geral, o valor $V$ a ser pago, em real, varia em função de:
\begin{itemize}
  \item tarifa base $F$: valor fixo, em real, cobrado no início da viagem;
  \item tempo $T$: tempo, em minutos, de duração da viagem;
  \item distância $D$: distância percorrida, em quilômetro.
\end{itemize}

Um desses aplicativos cobra R$ 2,00 de valor fixo, acrescido de R$ 0,26 por minuto de viagem e de R$ 1,40 por quilômetro rodado.

Nessas condições, a expressão que fornece o valor $V$ a ser pago por uma viagem desse aplicativo é:

\section{Solução Técnica}
A expressão correta para o valor da viagem é:
\[
V = 2,00 + 0,26T + 1,40D
\]
onde cada termo representa:
\begin{itemize}
  \item $2,00$: tarifa fixa;
  \item $0,26T$: custo pelo tempo de viagem;
  \item $1,40D$: custo pela distância percorrida.
\end{itemize}

\section{Explicação Didática}
O valor total é a soma da tarifa fixa com os custos proporcionais ao tempo e à distância. Todas as variáveis têm coeficientes numéricos que indicam o custo por unidade de tempo ou distância, respeitando o contexto da questão. É importante distinguir termos constantes (tarifa fixa) de termos variáveis multiplicados por coeficientes.

\section{Análise das Distratores}
Cada alternativa incorreta apresenta um tipo de erro conceitual:
\begin{itemize}
  \item Alternativa A: Confunde constante com variável multiplicada.
  \item Alternativa C: Omissão do custo por distância.
  \item Alternativa D: Esquecimento da tarifa fixa.
  \item Alternativa E: Ausência de valores monetários na fórmula.
\end{itemize}

\section{Perfil da Questão}
A questão trabalha a habilidade de construir e interpretar expressões algébricas simples em um contexto aplicado de matemática financeira. Possui baixa dificuldade e não requer cálculos complexos ou modelagem extensa.

\section{Conclusão}
A alternativa correta é a que apresenta a fórmula $V = 2,00 + 0,26T + 1,40D$.
Esta análise é consistente com o enunciado e as alternativas propostas.

\end{document}